\section{A worked capacity and decision example}\label{sec:C11:worked}
Consider an explicitly hypothetical decoder with 32 layers, 8 key-value heads, head dimension 128, and two-byte cached scalars. Equation~\eqref{eq:C11:kv} gives 131072 bytes per cached token. A sequence of 32768 tokens therefore requires 4 GiB of raw keys and values, while 131072 tokens require 16 GiB. These are arithmetic consequences of the stated configuration, not measurements of a named model or service. They exclude every other allocation and do not establish how many sessions fit on a particular accelerator. Replication, quantization, cache paging, and attention implementation can change the actual residency.

Suppose a deployment must hold ten such long sequences concurrently. The raw cache term alone becomes 160 GiB before weights or workspace. Moving half the sessions to host memory reduces device residency, but a resumed session may need cache transfer or recomputation. A throughput analysis that treats those sessions as permanently active can differ sharply from a workload in which users pause for minutes between tool decisions. Concurrency is therefore a property of arrival and activity patterns, not merely a maximum batch-size argument.

Now consider two histories $h_0$ and $h_1$ that are identical except for a valid user correction: one permits publishing an artifact and the other permits only preparing a draft. A summary function that omits the correction maps both histories to the same state $z$. Suppose the task requires a binary publish-or-hold action, the two histories are equally likely, and the correct actions differ. Any deterministic policy conditioned only on $z$ must be wrong for one history, giving error at least one half. Randomization cannot lower the average error because the optimal mixture still lacks the distinguishing bit. Unlimited additional reasoning on $z$ cannot recover an observation the representation discarded.

If the correction remains available in raw storage, retrieval can repair the failure, but only if the retrieval policy requests the right evidence before the consequential action. The memory system has then preserved capacity while the working-state policy has failed accessibility. If the correction was deleted from all retained evidence, retrieval cannot repair it. If the correction is present but labeled as an untrusted web instruction, its authority is misrepresented. These three cases can yield the same final wrong action, yet their engineering remedies are different.

This example also exposes a limitation of average semantic fidelity. The omitted correction can occupy a tiny fraction of a very long history while determining the entire action loss. A representation that preserves almost every token may be worse for the task than one that retains only a small structured set of obligations and evidence pointers. Token retention is not a monotone proxy for decision sufficiency. Conversely, a compact obligation list can fail when a later task needs an unforeseen detail, which is why retaining reversible access to raw history remains useful.

The proposed diagnostic is to create matched history pairs differing in one consequential fact and test whether a context policy preserves that distinction through length extension, compaction, interruption, and resumption. Positive results would support a specified invariance over the tested transformations. They would not prove that every unknown future distinction is preserved. This makes a carefully delimited behavioral claim possible without pretending that a generic summary is a sufficient statistic for all future tasks.

